

\subsection*{Theory}
The overarching concept of this project is ontological security, which is largely attributed to Giddens' \parencite*{giddensModernitySelfidentitySelf2009} \textit{Modernity and Self Identity}. Ontological security highlights the roll of everyday actions and their impact on self-identity and a sense of safety, thus creating a framework that relies on understanding the consistency of daily actions as well as the potential risks of these actions \Parencite[][p. 36]{giddensModernitySelfidentitySelf2009}. Trust constitutes one of the ways in which ontological security can be achieved through a emtotional and cognitive understanding of the world around someone; as Giddens \parencite*[][p. 38]{giddensModernitySelfidentitySelf2009} states, ``Basic trust is connected in an essential way to the interpersonal organisation of time and space,'' which relies on the actions and discourse that occur in interpersonal insteraction. For example, an interaction that leads to the intended outcome often stregnthens `basic-trust', whereas interactions that stray from the intended outcome often weaken it. Another crucial element to ontological security is routine, which of the daily actions that are repeated in ones environment that constitute a sense of being and create the framework for self-identity and ontological security \Parencite[][p. 39-40]{giddensModernitySelfidentitySelf2009}. The final of Giddens' \parencite*[][p. 43]{giddensModernitySelfidentitySelf2009} concepts discussed for the purposes of this research is social organisation. Social organisation comes from routines, trust, and mutual understanding, and this largely shapes one's self-identity in relation to the external world and other people \Parencite[][p. 43-44]{giddensModernitySelfidentitySelf2009}. Concepts like language also constitute social organisation as it requires a mutual understanding of what a word or an object is, what it does, and how it can be used, but through continual repetition of this language, it can form a trusted understanding that is understood by the self \Parencite[][p. 43-45]{giddensModernitySelfidentitySelf2009}. The combination of these concepts, both on an individual and societal level, create the basis of understanding ontological security, and if there is a disruption to any aspect of self-identity then ontological insecurity may arrise.

Expanding upon ontological security as proposed by Giddens \parencite*{giddensModernitySelfidentitySelf2009}, Mitzen \parencite*{mitzenOntologicalSecurityWorld2006} applies these concepts to the state by taking these concepts and applying them to the collective identities that constitute states. As Mitzen \parencite*[][p. 352]{mitzenOntologicalSecurityWorld2006} discusses, states seek ontological security as a result of the needs of their constituents who seek a stable identity and differentiate themselves from other societies. This need for stable identity by a state's constituents incentivizes maintaining and stregnthening national identity, which in many ways is driven by those who govern the state \Parencites[][p. 352]{mitzenOntologicalSecurityWorld2006}[][p. 29]{kocanIdentityOntologicalSecurity2023}. Furthermore, the actions of politicians can further show the effects of national idetities and how members of the collective idetity will act in ways that align with it despite possibly acting irrationally as a result \parencite[][p. 352]{mitzenOntologicalSecurityWorld2006}. Combining these concepts with theories such as securitisation theory, as Kočan (\cite*[][p. 28-30]{kocanIdentityOntologicalSecurity2023}; \cite[][p. 23-26]{buzan1998}) poses, provides a framework to analyse how the speech acts of politicians can shape the ontological (in)security. Discoursive actions by politicians lead to securitization, but this discourse both influences and is influenced by national identity and ontological security, further perpetuating a consisten narrative with the identity that has been formed \Parencites[][p. 29-31]{kocanIdentityOntologicalSecurity2023}[][p. 31-32]{buzan1998}. This ability of discourse to perpetuate identites, as well as the social power placed in the actions of politicians, provides the primary framework for utilizing political speeches to understand the ontological security of BiH \Parencite[][p. 32-33]{buzan1998}.

The combination of these theoretical concepts provides insight into the study of ontological security, identity, and securitisation at a state level, as carried out by politicians, but much of these studies focus on largely cohesive identities that underpin national identity. The case of BiH provides a unique approach to viewing how both political partisanship and ethnic ties affect and are affected by national identity and ontological security. As stated in the constitution of BiH \parencite*[][Annex 4]{oscemissiontobosniaandherzegovina1995}, all parlimentary, executive, and judicial institutions are divided evenly amongst the Bosniaks, Croats, and Serbs, thus leading to political institutions that are shaped by both political and ethnic lines. This disjunction allows for a nuanced approach to determining how BiH seeks and maintains its national identity and ontological security.

\subsection*{Methodology}

The methodology used for this project is a mixed methods apporach, involving both qualitative and quantitative content and discourse analysis of speeches, parlimentary debates, and party manifestos made by Bosnian politicians from 2010-2020. As of date, the dataset is comprised of 127,713 statements during parlimentary debates \Parencite[][]{mochtak2022}, 34 manifestos \Parencite{lfabilrrz25}, and 665 presidential speeches \Parencite{Speeches}. These texts will all be translated using OPUS-MT machine translation model designed for slavic languages \Parencite{tiedemann2020}. Although the translation of texts using machine translation is not as precise as human translated texts, due to the quantity of texts and methods to be used, this was determined to be the most effective method to translate texts in Bosnian, Croatian, and Serbian. This is is carried out as a means of input allignment for further analysis, as working with multiligual corpora provides significant complications on this scale \Parencite[][p. 9-10]{haukelicht2023}.

The methods used for discourse and content analysis are a mix of machine automated and manual techniques that provide a robust understanding of how political speeches in BiH are affected by identity and if they lead to ontological security or insecurity. Machine automated analysis will be carried out using multiple models, including a roBERTa based model trained using data from the Manifesto Project \parencite*{Burst:2024} to determine various political ideologies and perspectives. Sentiment analysis will use the CardiffNLP roBERTa based sentiment analysis model \parencite*{camacho-collados-etal-2022-tweetnlp} trained on twitter data, and topic modelling will be done through a version of BERTopic \parencite*{grootendorst2022bertopic} that will be fine-tuned using a training set of the full corpus. As well as the machine automated analysis, randomly sampled texts from the corpus will undergo manual discourse and content analysis to provide both validation for the use of automated methods and greater context to the use of political speech in BiH.